% LaTeX code starts here
% This document reconstructs a multi-page exam from provided images.
% Note: Page 4 of the original exam appears to be missing from the provided images.
% This reconstruction includes pages 1, 2, 3, and 5.

\documentclass{ctexart} % 使用 ctexart 文档类以支持中文

\usepackage[a4paper, margin=1in]{geometry} % 设置页面尺寸和边距
\usepackage{enumitem} % 用于自定义列表环境
\usepackage{array} % 用于增强 tabular 环境的功能
\usepackage{amsmath} % 用于数学公式
\usepackage{ragged2e} % 用于对齐控制
\usepackage{titlesec} % 用于自定义章节标题样式
\usepackage{pifont}

% Customization for different sections (optional, keeping it simple first)
% \titleformat{\section}{\Large\bfseries}{\thesection}{1em}{} % Example: section title large and bold

\begin{document}

% --- Page 1 ---
\pagestyle{plain} % Use plain page style for page numbers

% Header
\raggedright % Align left
课程编号: COM07053 \hfill 北京理工大学 2012-2013 学年第 2 学期

\vspace{1em} % Increase vertical space

\centering
\textbf{\fontsize{16pt}{18pt}\selectfont 2011 级数据库原理与设计期末试题 A 卷} % Exam Title

\vspace{1em}

% Student Info
\raggedright % Align left
班级\underline{\hspace{3cm}}\quad 学号\underline{\hspace{3cm}}\quad 姓名\underline{\hspace{3cm}}\quad 成绩\underline{\hspace{3cm}} % Student Info fields

\vspace{1em}

% Grade Summary Table
\centering
\begin{tabular}{|c|c|c|c|c|c|}
\hline
一 & 二 & 三 & 四 & 五 & 六 \\
\hline
 &  &  &  &  &  \\
\hline
\end{tabular}

\vspace{2em}

% Section One: Fill-in-the-blanks
\raggedright
\textbf{\large 一、 填空题 (每空 1 分, 共 20 分)}
\vspace{1em}

\begin{enumerate}[label=\arabic*.] % Use arabic numbers followed by a dot
    \item 一般来说, 数据模型包括三个方面, 分别是\underline{\hspace{2cm}}、\underline{\hspace{2cm}}和数据完整性约束。
    \item 数据库的数据独立性表现在两个方面, 分别是\underline{\hspace{2cm}}和\underline{\hspace{2cm}}。
    \item 在关系数据库中, 当数据更新操作出现破坏参照完整性而使数据不一致时, 系统可采用的策略有\underline{\hspace{2cm}}、\underline{\hspace{2cm}}和置空值。
    \item E-R 模型的主要概念是: 实体、\underline{\hspace{2cm}}和\underline{\hspace{2cm}}。
    \item 嵌入式 SQL 引入\underline{\hspace{2cm}}机制来协调面向集合和面向记录的不同处理方式。
    \item 在关系数据库的查询优化中, 可以把某些选择操作同它前面要执行的笛卡尔积结合起来, 合并成为一个\underline{\hspace{2cm}}操作。
    \item 在强制存取控制中, 仅当主体的许可证级别\underline{\hspace{2cm}}客体的密级时, 该主体才能读相应的客体。
    \item \underline{\hspace{2cm}}是数据库并发控制和恢复的基本单位。
    \item 为了提高多粒度封锁的检查效率, 在多粒度封锁机制中引入了\underline{\hspace{2cm}}锁。
    \item 数据库系统的三级模式结构分别是:\underline{\hspace{2cm}}、\underline{\hspace{2cm}}和\underline{\hspace{2cm}}。
    \item 恢复系统故障的方法是对于未完成的事务执行\underline{\hspace{2cm}}操作, 对于已完成的事务执行\underline{\hspace{2cm}}操作。
    \item 数据库系统遇到死锁时必须从死锁状态中恢复。解除死锁的方法是\underline{\hspace{4cm}}。
    \item 从低一级的范式通过模式分解达到高一级范式的过程称为\underline{\hspace{4cm}}。
\end{enumerate}

% Page 1 Footer
\vfill % Push content to the top
\centering
第 1 页, 共 5 页 % Assumed footer for page 1, matches overall count

\clearpage % End of Page 1

% --- Page 2 ---
\pagestyle{plain}

% Section Two: Judgement Questions
\raggedright
\textbf{\large 二、 判断题 (每小题 1 分, 共 10 分, 正确的填 T, 错误的填 F, 将结果填入下表)}
\vspace{1em}

% Answer Table for Single Choice
\centering
\begin{tabular}{|*{10}{c|}}
\hline
1 & 2 & 3 & 4 & 5 & 6 & 7 & 8 & 9 & 10 \\
\hline
 &  &  &  &  &  &  &  &  &  \\
\hline
\end{tabular}

\vspace{2em}

\begin{enumerate}[label=\arabic*.] % Use arabic numbers followed by a dot
    \item 概念模型一般是稳定的，只要用户的需求不变，概念模型也不变。
    \item 在关系数据库中，元组的各个分量是有序的，但是元组的次序是无关紧要的。
    \item 在关系运算中，只要参加运算的关系是有限的，结果关系也一定是有限的。
    \item 数据库的合法用户可以操纵数据库中的所有数据。
    \item 关系数据库中的视图只能建立在基表上。
    \item 同一关系内部属性之间不存在参照约束关系。
    \item 静态转储不影响事务的运行。
    \item 封锁粒度越大，需要的锁越多，系统开销也越大。
    \item 每个函数依赖集的最小函数依赖集都是唯一的。
    \item 只包含两个属性的关系一定属于 BCNF。
\end{enumerate}

\vspace{2em}

% Section Three: Single Choice Questions
\raggedright
\textbf{\large 三、 单项选择题 (每小题 1 分, 共 20 分)}
\vspace{1em}

% Answer Table for Single Choice
\centering
\begin{tabular}{|*{10}{c|}}
\hline
1 & 2 & 3 & 4 & 5 & 6 & 7 & 8 & 9 & 10 \\
\hline
 &  &  &  &  &  &  &  &  &  \\
\hline
11 & 12 & 13 & 14 & 15 & 16 & 17 & 18 & 19 & 20 \\
\hline
 &  &  &  &  &  &  &  &  &  \\
\hline
\end{tabular}

\vspace{2em}

% Single Choice Questions (Page 2 content)
\begin{enumerate}[label=\arabic*.]
    \item 负责数据库系统建立、维护和协调工作的人员是 ( ) \\
    A. 最终用户 \quad B. 应用程序员 \quad C. 系统分析员 \quad D. 数据库管理员
    \item 采用树型结构表示各类实体以及实体间的联系的模型是 ( ) \\
    A. 关系模型 \quad B. 层次模型 \quad C. 网状模型 \quad D. 对象模型
    \item 下列选项中，与视图对应的是 ( ) \\
    A. 外模式 \quad B. 模式 \quad C. 内模式 \quad D. 索引
    \item 关系中的列，称为 ( ) \\
    A. 元组 \quad B. 域 \quad C. 属性 \quad D. 目或度
    \item 若 R 有 m 个元组，S 有 n 个元组，则 R 和 S 的笛卡尔积有 ( ) 个元组 \\
    A. m \quad B. n \quad C. m+n \quad D. m$\times$n
    \item 在对两个关系 R 和 S 进行连接操作时候, 如果希望结果包含 R 和 S 中不满足连接条件的元组, 那么该执行的操作是 ( ) \\
    A. 内连接 \quad B. 左外连接 \quad C. 右外连接 \quad D. 完全外连接
    \item 对于小关系上的选择操作应该选用的实现方法是 ( ) \\
    A. 顺序扫描法 \quad B. 排序合并法 \quad C. 二分查找法 \quad D. 索引扫描法
\end{enumerate}

% Continue Section Three: Single Choice Questions (Page 3 content)
% Starting from question 8, continuing the list from previous page.
\begin{enumerate}[label=\arabic*., resume] % Continue the enumerate list
    \item 大型集中式数据库系统，查询优化的重点是使 ( ) 最小化 \\
    A. 访问存储器的代价 \quad B. 计算代价 \quad C. 内存使用代价 \quad D. 通信代价
    \item 在查询优化中应尽可能早地执行的操作是 ( )。\\
    A. 选择 \quad B. 连接 \quad C. 笛卡尔积 \quad D. 并
    \item 数据库管理系统能实现对数据库中数据的查询、插入、修改和删除等操作，这种功能称为 ( )。\\
    A. 数据定义功能 \quad B. 数据管理功能 \quad C. 数据操纵功能 \quad D. 数据控制功能
    \item 数据库系统中，用户和角色的关系是 ( ) \\
    A. 1:1 \quad B. 1:n \quad C. m:n \quad D. m:1
    \item 需要后备副本参与恢复的数据库故障是 ( ) \\
    A. 事务故障 \quad B. 系统故障 \quad C. 介质故障 \quad D. 事务回滚
    \item 下列数据完整性控制策略中数据控制能力最灵活的是 ( )。\\
    A. 默认值 \quad B. 规则 \quad C. 约束 \quad D. 触发器
    \item 事务并发操作不当可能导致一些问题, 例如丢失更新, 其原因在于并发操作破坏了事务的 ( ) \\
    A. 原子性 \quad B. 一致性 \quad C. 隔离性 \quad D. 持久性
    \item 能够保证可串行化的封锁协议是 ( ) \\
    A. 一级封锁协议 \quad B. 二级封锁协议 \quad C. 三级封锁协议 \quad D. 两阶段封锁协议
    \item 按照多粒度封锁协议, 如果某事务已加 SIX 锁, 则还可以继续施加的锁是 ( ) \\
    A. S \quad B. X \quad C. IS \quad D. IX
    \item 在函数依赖的范围内, 规范化程度最高的是 ( ) \\
    A. 2NF \quad B. 3NF \quad C. BCNF \quad D. 4NF
    \item 在分组查询中, 要去掉不满足条件的记录, 应当 ( ) \\
    A. 使用 WHERE 子句 \quad B. 使用 HAVING 子句 \\
    C. 先使用 HAVING 子句, 再使用 WHERE 子句 \quad D. 先使用 WHERE 子句, 再使用 HAVING 子句
    \item 在具有检查点的数据库系统中, 操作 \ding{172}写数据 \ding{173}写日志 \ding{174}写检查点记录 \ding{175}写重新开始文件 正确的顺序是 ( )。\\
    A. \ding{172}-\ding{173}-\ding{174}-\ding{175} \quad B. \ding{172}-\ding{175}-\ding{174}-\ding{173} \quad C. \ding{173}-\ding{174}-\ding{172}-\ding{175} \quad D. \ding{175}-\ding{174}-\ding{173}-\ding{172}
    \item 在关系数据库的逻辑结构设计中, 需要将一个 m:n 联系转换为一个关系模式。此时该关系模式的主键应该是 ( ) \\
    A. m 端实体的键 \quad B. n 端实体的键 \quad C. m 端和 n 端实体的键的组合 \quad D. 全键
\end{enumerate}

\vspace{2em}

% Section Four: Short Answer Questions
\raggedright
\textbf{\large 四、 简答题 (共 10 分, 请将答案写在答题纸上)} % Note: Original says 10 points total.
\vspace{1em}

\begin{enumerate}[label=\arabic*.]
    \item 简述数据库管理系统的主要功能。(4分)
    \item 简述聚集索引的概念。(2分)
    \item 按照软件工程的思想进行数据库设计需要经历哪几个阶段？(4分)
\end{enumerate}

\vspace{2em}

% Section Five: Comprehensive Questions (Partial, intro only on this page)
\raggedright
\textbf{\large 五、 综合题 (共 20 分, 请将答案写在答题纸上)}
\vspace{1em}

\begin{enumerate}[label=\arabic*.]
    \item 已知：关系模式 R(A, B, D, E, H)，R 上的函数依赖集 F=\{A$\rightarrow$D, AB$\rightarrow$DE, E$\rightarrow$H\}，请确定关系模式 R 的一个候选键，要求给出相关理由。(4分)
    % Note: Question 1 continues on the next page with sub-questions and more questions under Section Five.
\end{enumerate}


% Continue Section Five: Comprehensive Questions (Page 5 content)
% Resuming the enumeration for the Comprehensive Questions section.
\begin{enumerate}[label=\arabic*., resume] % Continue the enumerate list for Section Five
    \item % This is question 1 continued from page 3. The sub-questions follow.
    判断下列两个调度的是否冲突可串行化，并说明理由。(4分)
    S$_1$={$r_1$(x),$w_1$(x),$r_2$(y),$w_2$(y),$r_1$(y),$w_1$(y)} \\ % Using math environment for formulas
    S$_2$={$r_1$(x),$r_2$(y),$w_1$(x),$w_2$(y),$r_1$(y),$w_1$(y)} % Using math environment for formulas

    \item 证明：R 属于 3NF 但 R 属于 2NF。(3分) % Note: Seems like a typo in the question, should be "但不属于 BCNF" or similar. Reconstructing as written.

    \item 已知某关系数据库包含如下关系模式：
    图书 Book (编号 Bno, 书名 Bname, 出版社 Press, 定价 Price)
    读者 Reader (编号 Rno, 姓名 Rname, 读者类别 Type)
    读者类别信息 RType (读者类别 Type, 允许借阅的图书数目 Vmax)
    借阅信息 Borrow (图书编号 Bno, 读者编号 Rno, 借出日期 Ldate, 归还日期 Rdate)
    其中带下划线的属性为主键。未归还的图书归还日期默认为空 (NULL)。
    用 SQL 语言表示下列操作：（6分）
    \begin{enumerate}[label=(\arabic*)] % Use (1), (2), (3) style labels
        \item 将如下图书信息添加到数据库中: \\
        书名: Database, 编号: b010, 出版社: Elsevier, 定价: 99.
        \item 读者 r001 于 2013 年 5 月 19 日归还了编号为 b010 的图书。请据此修改数据库的相关内容。
        \item 统计每个读者未归还的图书数量。
    \end{enumerate}
    \begin{enumerate}[label=(\arabic*), resume] % Continue (4), (5), ... for this question
        \item 设计相应的触发器实现图书借阅规则：(4分) \\
        读者借阅的图书数量超过该读者可以借阅的图书数量时，拒绝借阅并给出提示信息。
    \end{enumerate}
\end{enumerate}

\vspace{2em} % Add space before the next section

% Section Six: Design Question
\raggedright
\textbf{\large 六、 设计题 (共 20 分)}
\vspace{1em}

\begin{enumerate}[label=\arabic*.]
    

    \item 某医院病房管理信息系统需要管理如下基本信息：一个科室有多个病房和多个医生，一个病房只能属于一个科室，一个医生也只属于一个科室，但负责多个病人的诊治并给出相应的诊治意见；每个病房有多张病床，一个病人只能占用一张病床；由于轮流值班的原因，一个病人可能会在不同的时间由不同的医生负责诊治。 % Introduction text for the design question, but it's under Section Five.
    % Based on the image, this seems to be part of the 종합题 section (五).
    % It then lists sub-questions (1), (2), (3). Let's label them as sub-items under item 3 of section 5.
    \begin{enumerate}[label=(\arabic*)] % Use (1), (2), (3) style labels
        \item 请上述信息对此管理系统进行概念建模, 用 E-R 图给出相应的概念模型。要求在 E-R 图中注明属性和联系的类型。(6分)
        % \begin{center}
        % [Placeholder for E-R Diagram] % Cannot generate image, add a placeholder comment
        % \end{center}
        \item 将所得的 E-R 模型转换成关系模型。(5分)
        \item 确定各个关系模式的函数依赖集，并指出各关系模式的主键。(5分)
        \item 判断各关系模式最高可达到第几范式？如果某个关系模式不是 3NF，请给出关系可能的 3NF 分解。(4分)
    \end{enumerate}
\end{enumerate}

% Page 5 Footer
\vfill
\centering
第 5 页, 共 5 页 % Matches the footer in the image

\end{document}
% LaTeX code ends here